% =============================================================================
% main.tex — research-gap-analysis presentation deck (Simple theme)
%
% Copy this file into:  ./slides/<topic-short-name>-gap-analysis/main.tex
% Also copy the build machinery from the paper-beamer-deck skill:
%   cp skills/paper-beamer-deck/templates/Makefile  ./slides/<topic>-gap-analysis/
%
% Compile with XeLaTeX (preferred):  make   (or xelatex twice)
%
% This deck summarizes a MULTI-PAPER analysis. It is structured around the
% analysis (landscape -> taxonomy -> comparison -> graded gaps -> directions),
% NOT around one paper. For all Beamer mechanics, layout robustness, and the
% mandatory PDF-visual-inspection discipline, follow paper-beamer-deck/SKILL.md.
%
% Path assumptions (same as a single-paper deck):
%   - Deck at  <workspace>/slides/<topic>-gap-analysis/
%   - Theme    <workspace>/beamerthemeSimple.sty  (no logos/assets)
%   - The reused Makefile sets TEXINPUTS=../../: so xelatex finds the theme .sty.
%   - Theme: loads <workspace>/theme.tex if present, else falls back to Simple.
%     Restyle all decks by editing <workspace>/theme.tex (see workspace AGENTS.md).
% =============================================================================
\documentclass[aspectratio=169]{beamer}
\InputIfFileExists{../../theme.tex}{}{\usetheme{Simple}}

\graphicspath{{figures/}{../../}}

\usepackage{booktabs}
\usepackage{amsmath}

% ---------- talk metadata (fill in) ----------
\title{Research Gap Analysis: <Topic>}
\subtitle{<One-sentence takeaway>}
\author{Presenter: <your name>}
\institute{N papers analyzed}
\date{\today}

\begin{document}

% ---------- 1. Title ----------
\begin{frame}[plain]
  \titlepage
\end{frame}

% ---------- 2. Outline ----------
\begin{frame}{Outline}
  \tableofcontents
\end{frame}

% =============================================================================
% Keep ONE main idea per slide. Put long reasoning in notes.md, not on slides.
% The comparison matrix is the #1 failure mode: SIMPLIFY or SPLIT it instead of
% shrinking the font. Compile + visually inspect after each section.
% =============================================================================

% ---------- 3. Scope & Executive Summary ----------
\section{Scope \& Summary}
\begin{frame}{Scope}
  \begin{itemize}
    \item Papers included / area analyzed
    \item Intentionally out of scope
  \end{itemize}
\end{frame}

\begin{frame}{Executive Summary}
  \begin{block}{Headline finding}
    <the single most important conclusion of the analysis>
  \end{block}
\end{frame}

% ---------- 4. Papers Analyzed ----------
\section{Papers Analyzed}
\begin{frame}{Papers Analyzed}
  \small
  \begin{enumerate}
    \item <Title> --- <Authors>, <Venue/Year>
    \item <Title> --- <Authors>, <Venue/Year>
    \item <Title> --- <Authors>, <Venue/Year>
  \end{enumerate}
\end{frame}

% ---------- 5. Problem Landscape ----------
\section{Problem Landscape}
\begin{frame}{Problem Landscape}
  % Prefer a diagram (problem-space map / layered system view) over text.
  \begin{itemize}
    \item The shared problem these papers address
    \item Why it matters
  \end{itemize}
\end{frame}

% ---------- 6. Taxonomy ----------
\section{Taxonomy}
\begin{frame}{Taxonomy}
  % A diagram or compact table organizing the papers by key dimensions.
  \begin{center}
  \begin{tabular}{@{}lll@{}}
    \toprule
    Paper & Dimension A & Dimension B \\
    \midrule
    P1 & \dots & \dots \\
    P2 & \dots & \dots \\
    \bottomrule
  \end{tabular}
  \end{center}
\end{frame}

% ---------- 7. Cross-Paper Comparison ----------
\section{Cross-Paper Comparison}
\begin{frame}{Cross-Paper Comparison (simplified)}
  % SIMPLIFIED matrix only. Full matrix stays in analysis/<topic>/paper-matrix.md.
  % Split across slides if it does not fit -- do not shrink to fit.
  \small
  \begin{center}
  \begin{tabular}{@{}llll@{}}
    \toprule
    Paper & Mechanism & Workload & Does NOT cover \\
    \midrule
    P1 & \dots & \dots & \dots \\
    P2 & \dots & \dots & \dots \\
    \bottomrule
  \end{tabular}
  \end{center}
\end{frame}

% ---------- 8. Coverage ----------
\section{Coverage}
\begin{frame}{What Is Well Covered}
  \begin{itemize}
    \item Mature problems / techniques across the set
  \end{itemize}
\end{frame}

\begin{frame}{What Is Missing}
  \begin{itemize}
    \item Missing cases / assumptions / evaluations / system settings
  \end{itemize}
\end{frame}

% ---------- 9. Candidate Gaps (one slide per strong gap) ----------
\section{Candidate Gaps}
\begin{frame}{Gap 1: <name> \hfill \textbf{[Strong]}}
  \begin{itemize}
    \item \textbf{Statement:} <gap>
    \item \textbf{Why it matters:} <\dots>
    \item \textbf{Why unsolved by the set:} <\dots>
  \end{itemize}
\end{frame}

\begin{frame}{Other Gaps (Medium / Weak)}
  \begin{itemize}
    \item \textbf{[Medium]} <gap> --- needs sharper framing / better evaluation
    \item \textbf{[Weak]} <gap> --- too broad / incremental / already addressed
  \end{itemize}
\end{frame}

% ---------- 10. Most Promising Directions ----------
\section{Directions}
\begin{frame}{Most Promising Directions}
  \begin{enumerate}
    \item <direction> --- what would make it publishable
    \item <direction>
  \end{enumerate}
\end{frame}

% ---------- 11. Weak / Dangerous Directions ----------
\begin{frame}{Weak / Dangerous Directions}
  \begin{itemize}
    \item <tempting but weak direction> --- why to avoid it
  \end{itemize}
\end{frame}

% ---------- 12. Recommended Next Papers ----------
\section{Next Steps}
\begin{frame}{Recommended Next Papers}
  \begin{itemize}
    \item <paper / lead> --- why relevant \hfill \emph{(to verify)}
  \end{itemize}
  \footnotesize Do not fabricate citations; mark unverified entries ``to verify''.
\end{frame}

% ---------- 13. Final Takeaway ----------
\begin{frame}{Final Takeaway}
  \begin{center}
    \Large <what the selected papers collectively imply>
  \end{center}
\end{frame}

\end{document}
